\scanpage{39}
% L1–2
\deleted{\deleted{\uline{Lemme \unsure{2|1}}} \inserted{\uline{Proposition}} Soit \deleted{\unsure{$F$|$J$}} $M$ un ensemble de morphismes \deleted{(\unsure{simples projectifs})} de préschémas loc. noethériens. On suppose que}

% L7–8, a); circled addition with pointer
\deleted{a) La projection $f:P^1_Y\to Y$ \deleted{est} \inserted{et $\mathrm{id}_Y:Y\to Y$ sont} $\in M$.}

% L9–10
\deleted{b) Si \deleted{\unsure{$f$}} $X\xrightarrow{f}Y\xrightarrow{g}Z$ $f,g\in M$, alors \deleted{\unsure{on est de}} $gf\in M$}

% L11–15
\deleted{c) Si $X=P^r_Y\xrightarrow{g}Y$ est le morphisme structural canonique, et \deleted{si} $X'=P^1_Y\times P^{r-1}_Y$ est muni \deleted{du} \deleted{\unsure{$X'$}} morphisme birationnel bien connu $X'\overset{f}{\unsure{\mathord{-}|\to}}X$, et si $gf\in M$, alors $g\in M$.}

% L5–7, inner relocation arrow to just above Sous ces conditions
\inserted{\deleted{d) Si $X\xrightarrow{f}Y\xrightarrow{g}Z$ \deleted{sont} tels que $f$ soit une immersion fermée, $g\deleted{f}\in M$ [et $gf$ simple] alors $gf\in M$}}

% L3–4, outer relocation arrow to just above Sous ces conditions
\inserted{\deleted{\unsure{e)|a)} Pour un morphisme $f:X\to Y$, l'appartenance à $M$ est locale sur $Y$}}

% L16–17
\deleted{Sous ces conditions, tout morphisme [simple] projectif $f:X\to Y$ est $\in M$.}

% L18–22
\deleted{En effet, on voit par récurrence sur $r$, utilisant a), b), c), que les $P^{\deleted{\unsure{n}}r}_Y\xrightarrow{t}Y$ sont $\in M$. Utilisant d), il s'ensuit que pour tt sous-préschéma fermé $X$ de $P^n_Y$ [simple$/Y$] la projection \deleted{sur} $Y$ est $\in M$.}

% archive sheet number: 38
